\paragraph*{04.03.04.01. Persönliche Ambitionen und Interessen Dritter und deren potenzielle Auswirkungen auf das Projekt beurteilen sowie diese Kenntnisse zum Nutzen des Projekts einsetzen}
\kcilabel{para:04.03.04.01_PersönlicheAmbitionenUndInteressenDritter}{04.03.04.01}
\label{para:04.03.04.01_PersönlicheAmbitionenUndInteressenDritter}

%%% Task
\par Als \gls{wsl} \gls{ricefw} identifizierte ich systematisch persönliche Ambitionen, Interessen und Einflussmechanismen aller Stakeholder (formale Rollen, informelle Netzwerke, Machtverhältnisse, Motivationen). Analysierte potenzielle Auswirkungen und nutzte dieses Wissen für Projekterfolg. Ziel: Verständnis der Interessenlage, frühzeitige Risikenerkennung, Chance-Nutzung.

%%% Action
\par Strukturierte, mehrstufige Stakeholderanalyse (formale Hierarchien, informelle Strukturen). Umfeldanalyse in \gls{iea} und \gls{ieb}, Auswertung von Projekt- und Organisationsdokumenten (\gls{r3}, \gls{lns}), bilaterale Gespräche mit Schlüsselpersonen (Geschäftsführung \gls{ais}, \gls{fnz}, Programmumfeld). Entwickelte Strategien zur Einbindung von Stakeholdern, steuerte Erwartungen bzgl. Ressourcenmangel (\gls{pb}, \gls{r3}/\gls{s4}, \gls{is}-Know-How).

%%% Result
\par Stakeholder-Analyse ermöglichte Betriebsverantwortlichen der \gls{ais} (\gls{iea}, \gls{ieb}) frühzeitig die Notwendigkeit der Ausschreibung des Plattformbetriebs an externe Partner zu verdeutlichen. Sicherte Unterstützung der Geschäftsführung.

%%% Reflect
\par Intensive und kontinuierliche Stakeholder-Analyse war entscheidend für Navigation komplexer Machtverhältnisse und Interessenkonflikte. Zukünftig: Praxis beibehalten, Erkenntnisse stärker in Projektkommunikation und Planung integrieren, proaktiv auf Konflikte reagieren.